%% kapa: exam=1 round=30 year=2019 period=1 form=A subject=real_estate subject_ko=부동산학원론 seq=28 number=108
%% source: Q-Net 감정평가사 기출문제 (cst00309) artlSeq=5205205 file=extracted/1008/5205205/2019년 제30회 감정평가사 1차시험 문제지/2019년 제30회 감정평가사 1차시험 문제지/2019년도 제30회 감정평가사 1차 1교시 A형.hwp
%% conversion: kordoc 4.24.51 -> tools/gen1.py (auto); review=proofread-sonnet; flags=none
\begin{question}{108}
A지역 임대아파트의 시장수요함수가 ${ Q_{ d } = 100- \frac{ 1 } { 2 } P }$ 이고, 시장공급함수는 ${ Q_{ s } = 20+ \frac{ 1 } { 3 } P }$ 이다. 정부가 임대료를 시장균형임대료에서 36만원을 낮추었을 경우 A지역 임대아파트의 초과수요량은? (단, ${ Q_{ d }}$: 수요량, ${ Q_{ s }}$: 공급량, ${ P }$: 임대료, 단위는 천호 및 만원이고, 다른 조건은 불변임)
\begin{choices}
  \item 30천호
  \item 32천호
  \item 40천호
  \item 52천호
  \item 70천호
\end{choices}
\end{question}
\answer{1}
